\documentclass[12pt, twoside]{article}
%\documentclass[12pt, twoside]{article}
\usepackage[letterpaper, margin=1in, headsep=0.2in]{geometry}
\setlength{\headheight}{0.6in}
%\usepackage[english]{babel}
\usepackage[utf8]{inputenc}
\usepackage{microtype}
\usepackage{amsmath}
\usepackage{amssymb}
%\usepackage{amsfonts}
\usepackage[nomessages]{fp} %\FPeval{\var-name}{2*sin(pi/6)}
\usepackage{siunitx} %units in math. eg 20\milli\meter
\usepackage{yhmath} % for arcs, overparenth command
\usepackage{tikz} %graphics
\usetikzlibrary{quotes, angles, arrows, arrows.meta}
\usepackage{pgfplots}
\usepgfplotslibrary{fillbetween}
\pgfplotsset{compat=1.16}
\usepackage{graphicx} %consider setting \graphicspath{{images/}}
\usepackage{parskip} %no paragraph indent
\usepackage{enumitem}
\usepackage{multicol}
\usepackage{venndiagram}

\usepackage{fancyhdr}
\pagestyle{fancy}
\fancyhf{}
\renewcommand{\headrulewidth}{0pt} % disable the underline of the header
\raggedbottom
\hfuzz=2mm %suppresses overfull box warnings

\usepackage{hyperref}

\fancyhead[LO]{La Scuola d'Italia / Dr.\ Huson / IB Math 11 (a) \\* 9 November 2026}
\fancyhead[RO]{\fbox{\begin{minipage}{6cm}\footnotesize Nickname:\\[0.5cm]\end{minipage}}}
\fancyfoot[C]{\thepage}

\begin{document}
%==================================== VERSION (a) ====================================
\subsubsection*{2.10 Classwork: Systems of Equations and Exponents (no calculator)}

\emph{Start in class and finish for homework. Write each step of your
working on its own line.}

\begin{enumerate}[itemsep=0.15cm]

%--- Problem 1: Two 2x2 systems by elimination, scaling both equations; first one scaffolded step by step [7 marks] ---
%--- Source: Original (AA SL 2.1, prior learning); 2-6Q review, parallel of 2-8CW P1(c) ---
\item[1.]
Solve each system of equations by elimination.

\begin{enumerate}[itemsep=0.15cm]
  \item $\begin{aligned}[t]
    4x + 3y &= 5 \qquad (1) \\
    6x - 5y &= 17 \qquad (2)
  \end{aligned}$
  \begin{enumerate}[label=(\roman*), itemsep=0.1cm]
    \item Multiply equation (1) by $3$ and equation (2) by $2$, so that
    both equations have $12x$. Write the two new equations. Multiply
    \emph{every} term. \vfill
    \item Subtract one new equation from the other to eliminate $x$.
    Solve for $y$. \vfill
    \item Substitute your value of $y$ into equation (1) to find $x$.
    \vfill
  \end{enumerate}
  \item $\begin{aligned}[t]
    5x + 2y &= 3 \\
    3x + 7y &= 25
  \end{aligned}$ \vfill
\end{enumerate}

\newpage

%--- Problem 2: Arithmetic sequence: u1 and d from two terms, sum of 20 terms; then p and q in Sn = pn^2 - qn from two sums; scaffolded [13 marks] ---
%--- Source: Original (AA SL 1.2, 2.1); modelled on the IB Paper 1 "find the parameters" pattern (u1 and d from two facts; p and q in Sn) ---
\item[2.]
\fbox{\begin{minipage}{0.9\textwidth}
\textbf{Finding two unknowns.} Each fact you are given becomes one equation.
Two facts give two equations, and you solve the system.
\end{minipage}}

\medskip
In an arithmetic sequence, $u_{4} = 17$ and $u_{10} = 41$.

\begin{enumerate}[itemsep=0.15cm]
  \item Use $u_{n} = u_{1} + (n - 1)d$ to write two equations in $u_{1}$
  and $d$. \vfill
  \item Solve your equations to find $u_{1}$ and $d$. \vfill
  \item Find the sum of the first $20$ terms. \vfill
\end{enumerate}

The sum of the first $n$ terms of a different arithmetic sequence is
$S_{n} = pn^{2} - qn$. You are given that $S_{2} = 2$ and $S_{5} = 35$.

\begin{enumerate}[itemsep=0.15cm]
  \item[(d)] Substitute $n = 2$ to show that $4p - 2q = 2$. \vfill
  \item[(e)] Use $S_{5} = 35$ to write a second equation in $p$ and $q$.
  \vfill
  \item[(f)] Solve your two equations to find $p$ and $q$. \vfill
  \item[(g)] $S_{1}$ is the sum of just the first term, so
  $u_{1} = S_{1}$. Find $u_{1}$. \vfill
\end{enumerate}

\newpage

%--- Problem 3: Curve y = A(2^x) + B through two points: two equations, A and B, evaluate, solve an exponential equation; scaffolded [9 marks] ---
%--- Source: Original (AA SL 2.1, 1.5); modelled on the IB Paper 1 "A and B of a curve through two points" pattern ---
\item[3.]
The curve $y = A\left(2^{x}\right) + B$ passes through the points
$(1,\ 7)$ and $(3,\ 19)$.

\begin{enumerate}[itemsep=0.15cm]
  \item Substitute $x = 1$, $y = 7$ to write an equation in $A$ and $B$.
  Do the same for the point $(3,\ 19)$. \vfill
  \item Solve the two equations to find $A$ and $B$. \vfill
  \item Find the value of $y$ when $x = 5$. \vfill
  \item Find the value of $x$ when $y = 35$. \vfill
\end{enumerate}

%--- Problem 4: Surds: find integers a, b from a product by equating rational and surd parts; scaffolded [5 marks] ---
%--- Source: Original (AA SL 1.7 prior learning); parallel of 2-6Q (b) P3(b) ---
\item[4.]
Find the integers $a$ and $b$ such that
$\left(3 + \sqrt{2}\right)\left(a + b\sqrt{2}\right) = 4 - \sqrt{2}$.

\begin{enumerate}[itemsep=0.15cm]
  \item Expand the left side. Collect it into a rational part and a
  $\sqrt{2}$ part. \vfill
  \item The rational parts on the two sides are equal, and the
  $\sqrt{2}$ parts are equal. Write two equations in $a$ and $b$, and
  solve them. \vfill
\end{enumerate}

\newpage

%--- Problem 5: Exponent techniques: exact powers table; power of a power; same-base equations; equation quadratic in 3^x with u = 3^x, scaffolded [21 marks] ---
%--- Source: Original (AA SL 1.5, 1.7, 2.10); 2-6Q review (A 3(a), B 3(a), B 3(c)) ---
\item[5.]
\begin{enumerate}[itemsep=0.15cm]
  \item Write each number as a power of $2$, $3$ or $5$.

  \smallskip
  {\renewcommand{\arraystretch}{1.8}
  \begin{tabular}{|c|c|c|c|c|c|c|c|c|c|}
  \hline
  $16$ & $32$ & $64$ & $27$ & $81$ & $25$ & $125$ & $\dfrac{1}{8}$ & $\dfrac{1}{9}$ & $\sqrt{5}$ \\
  \hline
  \hspace{0.9cm} & \hspace{0.9cm} & \hspace{0.9cm} & \hspace{0.9cm} & \hspace{0.9cm} & \hspace{0.9cm} & \hspace{0.9cm} & \hspace{0.9cm} & \hspace{0.9cm} & \hspace{0.9cm} \\[0.4cm]
  \hline
  \end{tabular}}

  \item Simplify. A power of a power multiplies the exponents:
  $\left(a^{m}\right)^{n} = a^{mn}$.

  \begin{tabular}{@{}p{0.25\textwidth}p{0.25\textwidth}p{0.33\textwidth}}
  (i) $\left(a^{3}\right)^{2}$ & (ii) $\left(2x^{4}\right)^{3}$ &
  (iii) $\left(3a^{2}b\right)^{2} \times 2a^{-3}b$
  \end{tabular}
  \vfill
  \item Solve each equation by writing both sides as powers of the same
  base.

  \begin{tabular}{@{}p{0.25\textwidth}p{0.25\textwidth}p{0.33\textwidth}}
  (i) $2^{x + 1} = 32$ & (ii) $9^{x} = 27^{x - 1}$ &
  (iii) $5^{2x} = \dfrac{1}{125}$
  \end{tabular}
  \vfill
  \item Consider the equation $3^{2x} - 10\left(3^{x}\right) + 9 = 0$.
  \begin{enumerate}[label=(\roman*), itemsep=0.1cm]
    \item Let $u = 3^{x}$. Explain why $3^{2x} = u^{2}$, and write the
    equation in terms of $u$. \vfill
    \item Solve the equation for $u$. \vfill
    \item Hence find the two values of $x$. \vfill
  \end{enumerate}
\end{enumerate}

\end{enumerate}

\newpage
%==================================== VERSION (b) ====================================
\fancyhead[LO]{La Scuola d'Italia / Dr.\ Huson / IB Math 11 (b) \\* 9 November 2026}
\subsubsection*{2.10 Classwork: Systems of Equations and Exponents (no calculator)}

\emph{Start in class and finish for homework. Write each step of your
working on its own line.}

\begin{enumerate}[itemsep=0.15cm]

%--- Problem 1: 2x2 by elimination scaling both; 3x3 by elimination with scaling (labelled equations) [11 marks] ---
%--- Source: Original (AA SL 2.1, systems of linear equations); 2-6Q (b) P2 review ---
\item[1.]
Solve each system of equations by elimination.

\begin{enumerate}[itemsep=0.15cm]
  \item $\begin{aligned}[t]
    5x + 2y &= 3 \\
    3x + 7y &= 25
  \end{aligned}$ \vfill
  \item $\begin{aligned}[t]
    3x + 2y - z &= 1 \qquad (1) \\
    2x - 3y + 4z &= 19 \qquad (2) \\
    4x + 5y + 2z &= 9 \qquad (3)
  \end{aligned}$

  \smallskip
  Number every new equation. Eliminate the \emph{same} variable twice,
  from two different pairs, to get equations (4) and (5) in two
  variables. Solve (4) and (5), then substitute back to find the third
  variable. \vfill \vfill
\end{enumerate}

\newpage

%--- Problem 2: Arithmetic sequence: u1 and d from a term and a sum; p and q in Sn = pn^2 - qn; u2 = S2 - S1; show un = 4n - 5 [18 marks] ---
%--- Source: Original (AA SL 1.2, 1.3, 2.1); modelled on the IB Paper 1 "find the parameters" pattern ---
\item[2.]
\fbox{\begin{minipage}{0.9\textwidth}
\textbf{Finding two unknowns.} Each fact you are given becomes one equation.
Two facts give two equations, and you solve the system.
\end{minipage}}

\medskip
In an arithmetic sequence, $u_{3} = 14$ and $S_{8} = 148$.

\begin{enumerate}[itemsep=0.15cm]
  \item Write two equations in $u_{1}$ and $d$. \vfill
  \item Find $u_{1}$ and $d$. \vfill
  \item Find $u_{20}$. \vfill
\end{enumerate}

The sum of the first $n$ terms of a different arithmetic sequence is
$S_{n} = pn^{2} - qn$. You are given that $S_{2} = 2$ and $S_{5} = 35$.

\begin{enumerate}[itemsep=0.15cm]
  \item[(d)] Find the value of $p$ and the value of $q$. \vfill
  \item[(e)] Find $u_{1}$ and $u_{2}$, and hence the common difference.
  \vfill
  \item[(f)] Show that $u_{n} = 4n - 5$. \vfill
\end{enumerate}

\newpage

%--- Problem 3: Curve y = A(2^x) + B through (-1, 4) and (2, 11); A and B; write as 2^(x+1) + 3; solve an exponential equation [9 marks] ---
%--- Source: Original (AA SL 2.1, 1.5); modelled on the IB Paper 1 "A and B of a curve through two points" pattern ---
\item[3.]
The curve $y = A\left(2^{x}\right) + B$ passes through the points
$(-1,\ 4)$ and $(2,\ 11)$.

\begin{enumerate}[itemsep=0.15cm]
  \item Find the value of $A$ and the value of $B$. \vfill
  \item Show that the equation of the curve can be written as
  $y = 2^{x + 1} + 3$. \vfill
  \item Find the value of $x$ when $y = 131$. \vfill
\end{enumerate}

%--- Problem 4: Surds: integers a, b from a product by equating parts; hence a quotient in the form p + q root 3 [5 marks] ---
%--- Source: Original (AA SL 1.7 prior learning); parallel of 2-6Q (b) P3(b) with an added "hence" part ---
\item[4.]
\begin{enumerate}[itemsep=0.15cm]
  \item Find the integers $a$ and $b$ such that
  $\left(2 + \sqrt{3}\right)\left(a + b\sqrt{3}\right) = 5 + 4\sqrt{3}$.
  \vfill
  \item Hence write $\dfrac{5 + 4\sqrt{3}}{2 + \sqrt{3}}$ in the form
  $p + q\sqrt{3}$, where $p, q \in \mathbb{Z}$. \vfill
\end{enumerate}

\newpage

%--- Problem 5: Exponent techniques: exact powers table (incl. fractional); power of a power with negative and fractional exponents; same-base equations; equation quadratic in 3^x [23 marks] ---
%--- Source: Original (AA SL 1.5, 1.7, 2.10); 2-6Q review (A 3(a), B 3(a), B 3(c)) ---
\item[5.]
\begin{enumerate}[itemsep=0.15cm]
  \item Write each number as a power of $2$, $3$ or $5$.

  \smallskip
  {\renewcommand{\arraystretch}{1.8}
  \begin{tabular}{|c|c|c|c|c|c|c|c|c|c|}
  \hline
  $32$ & $64$ & $81$ & $125$ & $\dfrac{1}{8}$ & $\dfrac{1}{9}$ & $\sqrt{5}$ & $4\sqrt{2}$ & $\dfrac{1}{5\sqrt{5}}$ & $\sqrt[3]{9}$ \\
  \hline
  \hspace{0.9cm} & \hspace{0.9cm} & \hspace{0.9cm} & \hspace{0.9cm} & \hspace{0.9cm} & \hspace{0.9cm} & \hspace{0.9cm} & \hspace{0.9cm} & \hspace{0.9cm} & \hspace{0.9cm} \\[0.4cm]
  \hline
  \end{tabular}}

  \item Simplify. Give your answers using positive exponents only.

  \begin{tabular}{@{}p{0.25\textwidth}p{0.33\textwidth}p{0.25\textwidth}}
  (i) $\left(2x^{4}\right)^{3}$ &
  (ii) $\left(3a^{2}b\right)^{2} \times 2a^{-3}b^{-4}$ &
  (iii) $\dfrac{\left(4a^{6}\right)^{\frac{3}{2}}}{\left(2a\right)^{3}}$
  \end{tabular}
  \vfill
  \item Solve each equation.

  \begin{tabular}{@{}p{0.25\textwidth}p{0.33\textwidth}p{0.25\textwidth}}
  (i) $25^{x + 1} = 125^{x - 1}$ & (ii) $4^{x} \times 8^{x - 1} = 32$ &
  (iii) $27^{x} = \dfrac{1}{9\sqrt{3}}$
  \end{tabular}
  \vfill
  \item By letting $u = 3^{x}$, solve the equation
  $9^{x} - 4\left(3^{x + 1}\right) + 27 = 0$. \vfill \vfill
\end{enumerate}

\end{enumerate}
\end{document}
